\documentclass{article}

\usepackage{makecell}

\title{Ch. 14 CW \#1}
\author{Principles of Microeconomics}
\date{\today}

\begin{document}

\maketitle

\subsection*{Instructions} Please answer the following on a separate sheet of paper and turn them in on Gradescope by next class.

\subsection*{Problems}

\begin{enumerate}

\item How and why does a firm's average-total-cost curve in the short run differ from its average-total-cost curve in the long run?

\item Define economies of scale and explain why they might arise. Define diseconomies of scale and explain why they might arise.

\item Consider the following cost information for a pizzeria:

	\begin{center}
	\begin{tabular}{c | c | c}
	\textbf{Quantity} & \textbf{Total Cost} & \textbf{Variable Cost} \\
	\hline
	0 dozen pizzas & \$300 & \$0 \\
	1 & 350 & 50 \\
	2 & 390 & 90 \\
	3 & 420 & 120 \\
	4 & 450 & 150 \\
	5 & 490 & 190 \\
	6 & 540 & 240 \\
	\end{tabular}
	\end{center}
	
	\begin{enumerate}
	
	\item What is the pizzeria's fixed cost?
	
	\item Construct a table in which you calculate the marginal cost per dozen pizzas using the information on total cost. Also, calculate the marginal cost per dozen pizzas using the information on variable cost. What is the relationship between these sets of numbers? Explain.
	
	\end{enumerate}
	
\item You cousin Vinnie owns a painting company with fixed costs of \$200 and the following schedule for variable costs:

	\begin{center}
	\begin{tabular}{| c | c | c | c | c | c | c | c |}
	\hline
	\makecell{\textbf{Quantity of}\\\textbf{Houses}\\\textbf{Painted per}\\\textbf{Month}} & 1 & 2 & 3 & 4 & 5 & 6 & 7 \\
	\hline
	\makecell{\textbf{Variable}\\\textbf{Costs}} & \$10 & \$20 & \$40 & \$80 & \$160 & \$320 & \$640 \\
	\hline
	\end{tabular}
	\end{center}
	
	Calculate average fixed cost, average variable cost, and average total cost for each quantity. What is the efficient scale of the painting company? 
	
\item The city government is considering two tax proposals:

	\begin{itemize}
	
	\item A lump-sum tax of \$300 on each producer of hamburgers
	
	\item A tax of \$1 per burger, paid by producers of hamburgers
	
	\end{itemize}
	
	\begin{enumerate}
	
	\item Which of the following curves -- average fixed cost, average variable cost, average total cost, and marginal cost -- would shift as a result of the lump-sum tax? Why? Show this in a graph. Label the graph as precisely as possible.
	
	\item Which of these same four curves would shift as a result of the per-burger tax? Why? Show this in a new graph. Label the graph as precisely as possible.
	
	\end{enumerate}
	
\item Jane's Juice Bar has the following cost schedules:

	\begin{center}
	\begin{tabular}{c | c | c}
	\textbf{Quantity} & \textbf{Variable Cost} & \textbf{Total Cost} \\
	\hline
	0 vats of juice & \$0 & \$30 \\
	1 & 10 & 40 \\
	2 & 25 & 55 \\
	3 & 45 & 75 \\
	4 & 70 & 100 \\
	5 & 100 & 130 \\
	6 & 135 & 165
	\end{tabular}
	\end{center}
	
	\begin{enumerate}
	
	\item Calculate average variable cost, average total cost, and marginal cost for each quantity.
	
	\item Graph all three curves. What is the relationship between the marginal-cost curve and the average-total-cost curve? Between the marginal-cost curve and the average-variable-cost curve? Explain.
	
	\end{enumerate}
	
\item Consider the following table of long-run total costs for three different firms:

	\begin{center}
	\begin{tabular}{| c | c | c | c | c | c | c | c |}
	\hline
	\textbf{Quantity} & 1 & 2 & 3 & 4 & 5 & 6 & 7 \\
	\hline
	\textbf{Firm A} & \$60 & \$70 & \$80 & \$90 & \$100 & \$110 & \$120 \\
	\hline
	\textbf{Firm B} & 11 & 24 & 39 & 56 & 75 & 96 & 119 \\
	\hline
	\textbf{Firm C} & 21 & 34 & 49 & 66 & 85 & 106 & 129 \\
	\hline
	\end{tabular}
	\end{center}
	
	Does each of these firms experience economies of scale or diseconomies of scale?

\end{enumerate}

\end{document}